\chapter{Low Self-Esteem (Loss of Confidence)}

\section*{Definition}
A negative perception of self-worth and abilities, often accompanied by self-doubt, social withdrawal, and difficulty asserting oneself in personal or professional settings.

\section*{Pathophysiology}
Self-esteem is influenced by prefrontal cortical regulation of amygdala and limbic system responses to social evaluation. Oestrogen modulates serotonin and dopamine neurotransmission; declining oestrogen reduces serotonergic tone, increasing sensitivity to negative social feedback and self-criticism. Physical symptoms of menopause (hot flushes, weight gain, skin/hair changes) may alter body image and self-perception. Sleep disruption and fatigue impair cognitive function, further undermining confidence in work and social roles.

\section*{Causes}
\begin{itemize}
  \item Perimenopausal and menopausal changes (physical and emotional)
  \item Major depressive disorder
  \item Generalised anxiety disorder
  \item Social anxiety disorder
  \item Body image disturbance (weight gain, skin/hair changes)
  \item Chronic illness or pain
  \item Workplace stress or discrimination
  \item Relationship difficulties or divorce
  \item History of trauma or abuse
  \item Perfectionism and high self-standards
\end{itemize}

\section*{When to See a Doctor}
See your doctor if:
\begin{itemize}
  \item Severe or worsening symptoms
  \item Low self-esteem with depressed mood, social withdrawal, work impairment, or thoughts of self-harm
  \item Accompanied by other concerning signs
\end{itemize}

\section*{Related Symptoms}
\begin{itemize}
  \item Depression
  \item Anxiety
  \item Mood swings
  \item Fatigue
  \item Body image concerns
\end{itemize}

\textbf{Important:} If this symptom persists, your doctor will check for serious underlying causes.

\section*{Self-Care / Home Management}
\begin{itemize}
  \item Rest and avoid activities that make it worse.
  \item Maintain adequate hydration and a balanced diet.
  \item Over-the-counter remedies may provide symptomatic relief (consult a pharmacist or doctor).
  \item Monitor symptoms and seek professional advice if they persist or worsen.
  \item Avoid self-medicating with prescription medications without medical guidance.
\end{itemize}

\section*{How Your Doctor Will Evaluate You}
\begin{itemize}
  \item A thorough history and physical examination are the first steps in evaluation.
  \item Laboratory tests (blood work, urinalysis) may be ordered based on clinical suspicion.
  \item Imaging studies (X-ray, ultrasound, CT, MRI) may be indicated when structural pathology is suspected.
  \item Specialist referral is arranged when the diagnosis remains uncertain or requires specific expertise.
\end{itemize}

\section*{Treatment Options}
\begin{itemize}
  \item Treatment targets the underlying cause identified through diagnostic evaluation.
  \item Supportive care and symptom management are provided while awaiting definitive therapy.
  \item Medications, lifestyle modifications, and physical therapies may be prescribed as appropriate.
  \item Follow-up ensures treatment effectiveness and allows timely adjustments to the care plan.
\end{itemize}

\clearpage